\subsection{Extension of the integral}

It follows from the definition of the space \(L_{F}^{p}\) that the subspace \(H_{F}\) of continuous functions with compact support is dense in \(L_{F}^{p}\) (§3, No.~4, Def. 2). Every continuous (for the topology of convergence in mean of order p) linear function, defined on \(H_{F}\) and taking its values in a complete Hausdorff topological vector space G, can therefore be extended by continuity in a unique manner, to a continuous linear function defined on \(L_{F}^{p}\) with values in G (GT, II, §3, No.~6, Th. 2 and III, §3, No.~1, Prop. 3).

Now, for every continuous function \(\mathbf{f}\) with compact support, with values in the Banach space \(\mathrm{F}\) , we have defined (in Ch. III, §3, No.~1) the integral \(\mu(\mathbf{f}) = \int \mathbf{f} d\mu\) with respect to \(\mu\) , which is an element of \(\mathrm{F}\) , and we have proved (Ch. III, §3, No.~2, Prop. 6) the inequality

\[
\left| \int \mathbf {f} d \mu \right| \leqslant \int | \mathbf {f} | d | \mu | = \mathrm{N} _ {1} (\mathbf {f}).\tag{1}
\]

This inequality proves that the linear mapping \(\mathbf{f} \mapsto \int \mathbf{f} d\mu\) of \(\mathcal{H}_{\mathrm{F}}\) into \(\mathrm{F}\) is continuous for the topology of convergence in mean in \(\mathcal{H}_{\mathrm{F}}\) . It can therefore be extended by continuity to the entire space \(L_{F}^{1}\) , and we may make the following definition:

DEFINITION 1. --- The functions belonging to \(\mathcal{L}_{\mathrm{F}}^{1}(\mathrm{X},\mu)\) are said to be integrable with respect to the measure \(\mu\) (or, again, that they are \(\mu\) -integrable). The integral (with respect to \(\mu\) ) of the integrable function \(\mathbf{f}\) is by definition the value at \(\mathbf{f}\) of the extension by continuity to \(\mathcal{L}_{\mathrm{F}}^{1}\) of the linear mapping \(\mathbf{g} \mapsto \int \mathbf{g} d\mu\) of \(\mathcal{K}_{\mathrm{F}}\) into \(\mathrm{F}\) ; it is again denoted \(\mu(\mathbf{f})\) or \(\int \mathbf{f} d\mu\) , or \(\int \mathbf{f}(x) d\mu(x)\) or \(\int \mathbf{f}\mu\) , or \(\int \mathbf{f}(x)\mu(x)\) .

Example. --- Let X be a discrete space, \(\mu\) a measure on X, and set \(\alpha(x) = \mu(\varphi_{\{x\}})\) for every \(x \in X\) . The functions in \(F_{F}^{1}\) are then integrable, in other words \(L_{F}^{1} = F_{F}^{1}\) ; moreover, for every function \(f \in L_{F}^{1}\) ,

\[
\int \mathbf {f} d \mu = \sum_ {x \in X} \alpha (x) \mathbf {f} (x).
\]

For, let \(\mathbf{f} \in \mathcal{F}_F^1\) ; we have \(|\mu|^*(|\mathbf{f}|) = \sum_{x \in X} |\alpha(x)| \cdot |\mathbf{f}(x)| < +\infty\) (§1, No.~3, Example); for every \(\varepsilon > 0\) , there exists a finite subset \(M\) of \(X\) such that

\[
\sum_ {x \in \mathrm{X} - \mathrm{M}} | \alpha (x) | \cdot | \mathbf {f} (x) | \leqslant \varepsilon .
\]

The function \(\mathbf{g}\) equal to \(\mathbf{f}\) at the points \(x\in M\) where \(|\mathbf{f}|\) is finite, and to 0 elsewhere, belongs to \(\mathcal{H}(\mathrm{X};\mathrm{F})\) and, by the conventions that have been made,

\[
| \mu | ^ {*} (| \mathbf {f} - \mathbf {g} |) \leqslant \sum_ {\boldsymbol {x} \in \mathrm{X} - \mathrm{M}} | \alpha (\boldsymbol {x}) | \cdot | \mathbf {f} (\boldsymbol {x}) | \leqslant \varepsilon ,
\]

which proves that \(\mathbf{f} \in \mathcal{L}_{\mathrm{F}}^{1}\) . On the other hand,

\[
\left| \mu (\mathbf {g}) - \sum_ {x \in X} \alpha (x) \mathbf {f} (x) \right| \leqslant \sum_ {x \in X - M} | \alpha (x) | \cdot | \mathbf {f} (x) | \leqslant \varepsilon ,
\]

whence the second assertion.

In other words, the \(\mu\) -integrable functions \(\mathbf{f}\) are those for which the family \((\alpha(x)\mathbf{f}(x))_{x\in X}\) is absolutely summable (GT, IX, §3, No.~6), and the integral \(\int \mathbf{f} d\mu\) is the sum of this family.

Since \(\mu(\mathbf{f})\) is continuous on \(\mathcal{L}_{\mathrm{F}}^{1}\) by definition, and since it takes its values in a Hausdorff space, we have \(\mu(\mathbf{f}) = 0\) for every function that belongs to the closure of 0 in \(\mathcal{L}_{\mathrm{F}}^{1}\) , that is, is negligible; if \(\mathbf{f}\) and \(\mathbf{g}\) are two equivalent integrable functions, then \(\mu(\mathbf{f}) = \mu(\mathbf{g})\) . In other words, the value of \(\mu(\mathbf{f})\) depends only on the class \(\widetilde{\mathbf{f}}\) of the integrable function \(\mathbf{f}\) ; it is again denoted \(\mu(\widetilde{\mathbf{f}})\) , and the function \(\widetilde{\mathbf{f}} \mapsto \mu(\widetilde{\mathbf{f}})\) is a continuous linear mapping of \(L_{F}^{1}\) into F. If a function f, with values in F and defined almost everywhere in X, is equivalent to an integrable function, we again say that f is integrable and we write \(\int f d\mu = \mu(\widetilde{\mathbf{f}})\) ; one defines similarly an integrable function with values in \(\overline{R}\) , defined and finite almost everywhere, as well as its integral.

